\titledquestion{Sorting Implementations and Analysis}

\begin{parts}

\part[3]
Assume all keys are distinct. In Bubble Sort with left-to-right passes, which of the following statements is true?

\begin{choices}
\choice In one pass, both the smallest and the largest key can move at most one position toward their final positions.
\choice In one pass, the largest key can move several positions to the right, while the smallest key can move at most one position to the left.
\choice In one pass, the smallest key can move several positions to the left, while the largest key can move at most one position to the right.
\choice In one pass, both the smallest and the largest key can move several positions toward their final positions.
\end{choices}

\begin{solution}
B
\end{solution}

\part[4]
Complete the following Insertion Sort implementation for a singly linked list. The function sorts the list in increasing order by relinking the existing nodes and stores the number of inversions in \texttt{*inversions}. The input list contains distinct values.

Each blank should contain exactly one statement or expression.

\begin{lstlisting}
struct ListNode {
    int val;
    ListNode* next;
    ListNode(int x) : val(x), next(nullptr) {}
};

ListNode* insertionSortAndCount(ListNode* head, long long* inversions) {
    ListNode dummy(0);
    int seen = 0;
    *inversions = 0;

    while (head != nullptr) {
        ListNode* node = head;
        head = head->next;

        ListNode* prev = &dummy;
        int before = 0;

        while (________________________________________) { // (1)
            prev = prev->next;
            ++before;
        }

        ________________________________;                 // (2)
        ________________________________;                 // (3)
        *inversions += __________________;                // (4)
        ++seen;
    }

    return dummy.next;
}
\end{lstlisting}

\begin{solution}

\begin{enumerate}

\item[(1)] prev->next != nullptr \&\& prev->next->val < node->val

\item[(2)] node->next = prev->next

\item[(3)] prev->next = node

\item[(4)] seen - before

\end{enumerate}
\end{solution}

\part[2]
Consider the following swapping implementation of Insertion Sort. Immediately before each execution of \texttt{std::swap}, which of the following statements are always true? Select all that apply.

\begin{lstlisting}
for (int i = 1; i < n; ++i) {
    for (int j = i; j > 0 && a[j - 1] > a[j]; --j) {
        // assertion here
        std::swap(a[j - 1], a[j]);
    }
}
\end{lstlisting}

\begin{choices}
\choice The value originally at index \(i\) is currently at index \(j\).
\choice The subarray \(a[0..i]\) is already sorted.
\choice The value at \(a[j]\) is smaller than every element in \(a[0..j-1]\).
\choice The entries \(a[j+1..i]\), if any, are in nondecreasing order and are all larger than \(a[j]\).
\end{choices}

\begin{solution}
AD
\end{solution}

\part[2]
What is the minimum number of pairwise comparisons required to find the smallest of \(n\) items?

\begin{solution}

$n-1$ 次。

\end{solution}

\end{parts}
